
%%%%%%%%%%%%%%%%%%%%%%%%%%%%%%%%%%%%%%%%%%%%%%%%%%%%%%%%%%%%%%%%%%%%%%%%%%%%%%
% Detta är ett exempel på ett latexdokument.
% 
% Alla dokument består av följande delar:
%
%          \documentclass[optioner]{dokumentklass}
%            ...inställningar...
%          \begin{document}
%            ...text...
%          \end{document}
%
% Som ni kanske redan har förstått är används procent (%) för
% kommentarer.
%%%%%%%%%%%%%%%%%%%%%%%%%%%%%%%%%%%%%%%%%%%%%%%%%%%%%%%%%%%%%%%%%%%%%%%%%%%%%%

\documentclass[titlepage,a4paper]{article}
\usepackage[pdftex]{graphicx}
\usepackage{anysize}
\marginsize{3cm}{3cm}{3cm}{3cm}
\usepackage[T1]{fontenc}                % För svenska bokstäver
\usepackage[swedish]{babel}      
\usepackage{epstopdf}     % För svensk avstavning och svenska
\usepackage[osf]{mathpazo} % Palatino with smallcaps and oldstyle numbers
\usepackage[scaled]{helvet} % Helvetica, scaled 95%
                  

\begin{document}

\begin{titlepage}
\begin{center}
\includegraphics[width=0.30\textwidth]{./logo}~\\[1cm]
\textsc{\LARGE SRS}\\
\textup{Programvaruutv.}\\[1cm]
Martin Lichota\\


\end{center}
\end{titlepage}
                   % Skriver ut rubriken som vi
                                % deklarerade ovan med \title, \author
                                % och eventuellt \date

\section{Inledning}          % Detta kommando gör en rubrik

Syftet med laborationen var att studera optokomponenter i form av dioder. Laborationen bestod av tre olika moment där dioders absorption och emission mättes, diod som solcell studerades samt en diod tillverkades.


\section{Utförande}             % Detta blir en underrubrik
\subsection{Lysdiod - emission och absorption}
Med hjälp av en monokromator, i vilken ljuset delades upp med ett glasprisma, belystes en diod. Dioden reagerade på en specifik våglängd som gav upphov till en ström. Detta kunde registreras på datorn.

Genom att istället låta dioden belysa en fotodetektor kunde emissionen mätas. Även detta registrerades på datorn.

\subsection{Solcellen}
En diod kopplades upp enligt kretsschemat i laborationshandledningen. Genom att variera belastningsresistansen erhölls olika ström- respektive spänningsvärden vid konstant belysning(på dioden/solcellen). 
Dessa värden användes sedan för att plotta ett I-U - diagram. Även den maximala effekten beräknades.

\subsection{Tillverkning och analys av en GaP-diod}
En diod skulle tillverkas genom att legera GaP med tenn och en zinklegerad guldplatta. Därefter studerades diodens karakteristik; diodens färg(som den emitterade) samt diodens absorptionsförmåga.

%-------------------------------------------------------------------
\pagebreak
\section{Analys och bearbetning av experimentell data}
%-------------------------------------------------------------------
\subsection{Lysdiod - emission och absorption}
\begin{figure}[htb]
\centering
\includegraphics[scale=.6]{osynligt.eps}
\caption{Absorption (rött) och emission (blått) som funktion av våglängd (nm) - Osynligt ljus.}
\end{figure} 

\begin{figure}[htb]
\centering
\includegraphics[scale=.6]{synligt.eps}
\caption{Emission (blått) som funktion av våglängd (nm) - Synligt ljus.}
\end{figure} 

Det synliga ljuset var blått och har våglängden 
\begin{math} \lambda = 455*10^{-9} m.
\end{math}

Ur formeln 

\begin{equation}
E_{g}= \frac {hc}{\lambda} \label{ekv1} 
\end{equation}

fås $E_{g}=2,7eV$.

Det osynliga ljusets våglängd är enligt mätvärden ca $950*10^{-9}m$.
På samma sätt som ovan fås motsvarande bandgapet till $E_{g}=1,3 eV$.

%-------------------------------------------------------------------

\subsection{Solcellen}
\begin{figure}[htb]
\centering
\includegraphics[scale=.6]{solcell1.eps}
\caption{Ström (mA) som funktion av spänning (V). Kurvan representerar solcellens karakteristik.}
\end{figure} 
\begin{figure}[htb]
\centering
\includegraphics[scale=.6]{solcell2.eps}
\caption{Effekt (mW) som funktion av resistans ($k\Omega)$. Kurvan representerar solcellens karakteristik.}
\end{figure} 
Ur figur 4 avläses maxeffekten till 8,6 mW.

%--------------------------------------------------------------
\pagebreak
\subsection{Tillverkning och analys av en GaP-diod}
Ström utan belysning: 3,5 mA.\\
Ström vid belysning; 4,5 mA.\\
Emitterat ljus: Orange.
\begin{figure}[hb]
\centering
\includegraphics[scale=.5]{opt-GAP.png}
\caption{Diodens karakteristik, ström (mA) som funktion av spänning (V). Pilen till vänster har värdet -4 V och pilen till höger cirka 0,4 V.}
\end{figure} 



\section{Sammanfattning}
Dioders karakteristik stämde till stor del överens med teorin. I det första momentet blev det som förväntat emission och absorption inom en viss våglängd för en given diod. Nyckeln till fenomenen ligger i ämnets bandgap, med vars hjälp man kan identifiera det. Dioden med det osynliga ljuset är gjord av kisel (bandgap på 1,11 eV vid 302K) och den som ger synligt ljus är gjord av zinkselenid (bandgap på 2,7 eV vid 302K).

Från det andra momentet kan det konstateras att man med en lämplig belastning kan utvinna ström och därmed tillverka en solcell. Detta sker då man teoretiskt sätt befinner sig i den fjärde kvadranten vilket medför till negativ ström. Den största arean mellan kurvan och axlarna (produkten av ström och spänning) ger den maximala effekten 8,6 mW.

I det tredje momentet inträffade ett intressant fenomen vid negativ spänning, kallad \textit{break down}. I övrigt stämde diodens karakteristik överens med det förväntade, nämligen tröskelspänning vid 0,6 V (0,4 V på laborationen). Dioden påvisade ökad ström vid belysning och fungerade därför som solcell; den emitterade även orange ljus.

\end{document}                 % The input file ends with this command.